\section{What remains: the threshold interval $[2,\tfrac{56501}{6400}]$}\label{sec:remains}

By Theorem~\ref{thm:h883} and Proposition~\ref{prop:barrier}, the set of admissible thresholds is a closed half-line
$[c^*,\infty)$ with
\[
2\ \le\ c^*\ \le\ \tfrac{56501}{6400}=8.828\ldots
\]
(the set is upward closed, and closed because for each instance both sides depend continuously on $c$). This section
records what is known about the interior of that interval: an exact truncation available at \emph{every} threshold, the
inequality with threshold $2$ for systems with few primes and for systems of large mean, and the precise statement of
what is left. Nothing here narrows the interval.

\begin{remark}[the barrier family is not tight at $2$]\label{rem:notight}
On the family of Proposition~\ref{prop:barrier} the two sides are $\binom r2(2-c)$ and $r-1$, so equality holds at
$c=2-\tfrac2r$ and at $c=2$ the left side vanishes while the right side is $r-1$. The family therefore forces
$c^*\ge2$ only in the limit $r\to\infty$, and provides no evidence whatever that $c^*=2$; the true threshold may lie
anywhere in $[2,\tfrac{56501}{6400}]$.
\end{remark}

\begin{lemma}[primorial truncation]\label{lem:primorial}
Let $d\ge0$ be an integer and let $Q_d=p_1\cdots p_d$ be the $d$-th primorial ($Q_0=1$, $Q_1=2$, $Q_2=6$, $Q_3=30$,
$Q_4=210$, $Q_5=2310$, $Q_6=30030$, $Q_7=510510$). Given an atom system $S$ and $m\ge1$, let $S_0$ consist of the atoms
$p^j$ with $p^jQ_d\le m$. Then for every $k\le m$ and every real $c\ge d$,
\[
\bigl(S(k)-c\bigr)^+=\bigl(S_0(k)-c\bigr)^+ ,
\]
so $\sum_{k\le m}(S(k)-c)^+=\sum_{k\le m}(S_0(k)-c)^+$; and $S_0\le S$ gives
$\sum_{b\in I}(S_0(b)-1)^+\le\sum_{b\in I}(S(b)-1)^+$ for every window. In particular a counterexample to $(H_c)$
yields one supported on $p^j\le m/Q_d$. The primorial is optimal: for every $Q'>Q_d$ there are $S$ and $m$ for which
the identity fails at $c=d$.
\end{lemma}
\begin{proof}
Suppose an omitted atom $q=p^j$ (that is, $qQ_d>m$) divides $k\le m$, and write $k=qt$. Then $t\le m/q<Q_d$, so
$\omega(t)\le d-1$: an integer with $d$ distinct prime factors is at least the product of the $d$ smallest primes,
namely $Q_d$. Since $q$ is a prime \emph{power}, $\omega(k)\le\omega(t)+1\le d$. The per-prime capacity now gives
\[
S(k)=\sum_{p\mid k}\ \sum_j\alpha_{p,j}[p^j\mid k]\ \le\ \sum_{p\mid k}1=\omega(k)\le d\le c ,
\]
and $S_0\le S$, so both hinges vanish at $k$. At every other $k\le m$ no omitted atom of positive weight divides $k$,
so $S(k)=S_0(k)$. For optimality take $q$ the least prime exceeding $Q_d$, $m=Q_dq$, and unit weights on
$p_1,\dots,p_d$ and $q$: then $S(m)=d+1$ and the hinge at $k=m$ equals $1$, while any cutoff $m/Q'$ with $Q'>Q_d$
omits $q$ and the hinge collapses to $0$.
\end{proof}

The truncation of Lemma~\ref{lem:h97-2} is the case $d=7$ with $2^{16}$ in place of $Q_7=510510$, and is weaker
($2^{16}<Q_7$); we keep it in Section~\ref{sec:h97} because that is the form which has been formalised.
The point of Lemma~\ref{lem:primorial} is that it is available at \emph{every} threshold: at $c=2$ one may truncate at
$p^j\le m/6$, whereas no argument of the $2^{16}$ type truncates at all below $c=7$. It is a prerequisite for the two
propositions below and for any attack on $c<7$. Its price is that the floor
$\lfloor m/q\rfloor\ge\tfrac{Q_d}{Q_d+1}\,\tfrac mq$ degrades from $1-2^{-16}$ to $\tfrac67$ at $d=2$, and that the
truncated system may still carry every prime power up to $m/6$.

\begin{proposition}[threshold $2$ for at most four primes]\label{prop:fourprimes}
Let $S$ be an atom system in which at most four primes carry positive weight. Then
$\sum_{k\le m}(S(k)-2)^+\le\sum_{b\in I}(S(b)-1)^+$ for every $m\ge1$ and every window $I$ of $m$ consecutive positive
integers.
\end{proposition}
\begin{proof}
Padding with zero functions, write $S=f_1+f_2+f_3+f_4$ with $f_i(n)=\sum_j\alpha_{p_i,j}[p_i^j\mid n]\in[0,1]$, put
$g_{ij}=(f_i+f_j-1)^+$ and $G=\tfrac13\sum_{i<j}g_{ij}$. Each $f_i$ occurs in exactly three of the six pairs, so
$\sum_{i<j}g_{ij}\ge\sum_{i<j}(f_i+f_j-1)=3S-6$; the sum is also nonnegative, whence $G\ge(S-2)^+$. For the other
half, the six pairs split into the three perfect matchings $\{12,34\}$, $\{13,24\}$, $\{14,23\}$, and
$(a-1)^++(b-1)^+\le(a+b-1)^+$ for $a,b\ge0$; applying this to each matching and averaging gives $G\le(S-1)^+$.

Fix a pair $p,q$ and let $A_a=\sum_{j\le a}\alpha_{p,j}$, $B_b=\sum_{j\le b}\alpha_{q,j}$ be the cumulative weights, so
that $f_p(n)=A_{\min(v_p(n),J_p)}$ where $J_p$ is the largest exponent occurring at $p$. With $\varphi(u)=(u-1)^+$ put
\[
\beta_{a,b}=\varphi(A_a+B_b)-\varphi(A_{a-1}+B_b)-\varphi(A_a+B_{b-1})+\varphi(A_{a-1}+B_{b-1}).
\]
Then $\beta_{a,b}\ge0$: with $\delta=A_a-A_{a-1}\ge0$ the function $u\mapsto\varphi(u+\delta)-\varphi(u)$ is
nondecreasing because $\varphi$ is convex, and $B_b\ge B_{b-1}$. Since $A_a\le1$ and $B_b\le1$, the terms with index
$0$ vanish, and telescoping in both variables gives the finite expansion
$g_{pq}(n)=\sum_{a,b\ge1}\beta_{a,b}[p^aq^b\mid n]$. Every modulus $d$ satisfies
$\#\{b\in I:d\mid b\}\ge\lfloor m/d\rfloor$, so $\sum_{k\le m}g_{pq}(k)\le\sum_{b\in I}g_{pq}(b)$. Summing the six
pairs,
\[
\sum_{k\le m}(S(k)-2)^+\le\sum_{k\le m}G(k)\le\sum_{b\in I}G(b)\le\sum_{b\in I}(S(b)-1)^+ . \qedhere
\]
\end{proof}

The certificate $G$ depends only on the atom system, not on $m$ or on the window. It does not survive a fifth prime:
for five weights equal to $\tfrac12$ one has $(S-2)^+=\tfrac12$ while every pair term vanishes.

\begin{proposition}[threshold $2$ for systems of large mean]\label{prop:densetwo}
Let $S$ be an atom system, $m\ge1$, and let $S_0$ be its truncation at $p^jQ_d\le m$ as in
Lemma~\ref{lem:primorial}. Write $H(f)=\sum_{p,j}\beta_{p,j}/p^j$ for the mean of a system with nonnegative
prime-power increments $\beta_{p,j}$. If, for an integer $c\ge2$ and a real $h>0$,
\[
H(S_0)\ge h\qquad\text{and}\qquad \frac{Q_c}{Q_c+1}\,h-\frac{h^c}{(c+1)!}\ \ge\ 1 ,
\]
then $(H_c)$ holds for $S$, $m$ and every window. In particular, taking $c=2$, $Q_2=6$ and
$h=h_0:=\tfrac{18-\sqrt{30}}7=1.78896\ldots$, every system whose $m/6$-truncation has mean at least $h_0$ satisfies
the hinge inequality with threshold $2$; the rational value $\tfrac{179}{100}$ suffices, since
$\tfrac67\cdot\tfrac{179}{100}-\tfrac16\bigl(\tfrac{179}{100}\bigr)^2-1=\tfrac{113}{420000}>0$.
\end{proposition}
\begin{proof}
Put $V=(h/H(S_0))\,S_0$, so that $V\le S_0\le S$, the per-prime capacities are still at most $1$, and $H(V)=h$
exactly; note that the hypothesis is used only through this rescaling, so no upper bound on $H(S_0)$ is needed. For
$q\le m/Q_c$ we have $\ell=\lfloor m/q\rfloor\ge Q_c$ and $m/q<\ell+1$, hence
$\lfloor m/q\rfloor\ge\tfrac{Q_c}{Q_c+1}\tfrac mq$ and $\sum_{k\le m}V(k)\ge\tfrac{Q_c}{Q_c+1}mh$. For $u_i\in[0,1]$
and an integer $c\ge2$ one has $(\sum_iu_i-c)^+\le e_c(u)/(c+1)$ (the difference is minimised at a vertex of the cube
after the usual compression, where it reduces to $\binom kc\ge(c+1)(k-c)^+$), so Lemma~\ref{lem:h97-1} gives
$\sum_{k\le m}(V(k)-c)^+\le mh^c/(c+1)!$. Subtracting, $\sum_{k\le m}\min(V(k),c)\ge m$, and $V\le S$ gives
$\sum_{k\le m}\min(S(k),c)\ge m$. Now the affine certificate $F=S-1$ is pointwise at most $(S-1)^+$, and
\[
\sum_{b\in I}F(b)=\sum_{p,j}\alpha_{p,j}\#\{b\in I:p^j\mid b\}-m\ \ge\ \sum_{k\le m}S(k)-m
\ \ge\ \sum_{k\le m}S(k)-\sum_{k\le m}\min(S(k),c)=\sum_{k\le m}(S(k)-c)^+ ,
\]
while $\sum_{b\in I}F(b)\le\sum_{b\in I}(S(b)-1)^+$.
\end{proof}

The same argument with $V=\tau S$ ($0\le\tau\le1$, no truncation) shows that
$\mu-\tfrac16H(S)^2\ge1$ implies $(H_2)$, where $\mu=\tfrac1m\sum_{k\le m}S(k)$; and under the hypotheses below the
choice $\tau=1$ is optimal.

\begin{proposition}[the remaining case]\label{prop:remains}
The hinge inequality with threshold $2$ holds for all atom systems, all $m$ and all windows if and only if it holds
for every $S$, $m$, $x$ subject to: all atoms satisfy $p^j\le m/6$; at least five primes carry positive weight;
$m\ge66$; $\sum_{k\le m}(S(k)-2)^+>0$; $H(S)<h_0=\tfrac{18-\sqrt{30}}7$; and $\mu-\tfrac16H(S)^2<1$ with
$\mu=\tfrac1m\sum_{k\le m}S(k)$.
\end{proposition}
\begin{proof}
One direction is trivial. Conversely, let $(S,m,x)$ violate $(H_2)$ and replace $S$ by its $m/6$ truncation, which by
Lemma~\ref{lem:primorial} leaves the left side unchanged and does not increase the right side. At most four active
primes is excluded by Proposition~\ref{prop:fourprimes}. A positive left side forces some $k\le m$ with $S(k)>2$,
hence $\omega(k)\ge3$; with five active primes all of whose atoms are at most $m/6$, the fifth smallest active prime
is at least $11$, so $m\ge66$. Mean at least $h_0$ is excluded by Proposition~\ref{prop:densetwo} applied with
$h=h_0$, and $\mu-\tfrac16H^2\ge1$ by the remark following it.
\end{proof}

The region described by Proposition~\ref{prop:remains} is not empty --- for instance $m=210$ with the atoms
$2,4,3,13,19,23,29$ of weights $\tfrac23,\tfrac13,1,1,1,\tfrac23,1$ satisfies every hypothesis --- so the
proposition is a genuine reduction and not a disguised proof. Since truncation gives $\mu\ge\tfrac67H$, the last two
hypotheses together already force $H<h_0$ or $H>h_1=\tfrac{18+\sqrt{30}}7$, the second being excluded only by the
rescaling in the proof of Proposition~\ref{prop:densetwo}. Analogous reductions hold at every integer threshold
$c\ge3$, with $Q_c$ in place of $6$ and $m\ge Q_{c+1}$, but without the five-prime clause: no analogue of
Proposition~\ref{prop:fourprimes} is known for $c\ge3$.

\subsection*{Beyond four primes}

Proposition~\ref{prop:fourprimes} is the base case of an induction: an atom system is finitely
supported, so $(H_2)$ holds in general as soon as it holds for every bounded number of active
primes. The pair certificate itself does not survive a fifth prime, but two different mechanisms do,
and a third is provably unavailable.

\begin{lemma}[Boolean antichain bound]\label{lem:antichain}
Let $2\le a<b<c<d$ be pairwise coprime integers and $Q\ge2$ an integer. Let $\mathcal M_Q$ be the
family of inclusion-minimal nonempty $A\subseteq\{a,b,c,d\}$ with $\prod A/\max A\ge Q$. Then
$Q\sum_{A\in\mathcal M_Q}\bigl(\prod A\bigr)^{-1}\le1$, with maximum $\tfrac{101}{105}$ attained at
$(\{2,3,5,7\},Q=2)$.
\end{lemma}
\begin{proof}
A pair is feasible exactly when its smaller member is $\ge Q$; a triple exactly when the product of
its two smaller members is $\ge Q$, and is minimal exactly when it contains no feasible pair; the
quadruple is feasible exactly when $abc\ge Q$ and minimal exactly when $bc<Q$. This gives nine
cases according to the position of $Q$ relative to $a,b,c,ab,ac,bc,abc$, and in each the displayed
bound is $\le1$: all rows but the first are bounded by $\tfrac1c+\tfrac3d\le\tfrac4c\le\tfrac45$
(pairwise coprimality forces $c\ge5$), and the first row is $a\,e_2(\tfrac1a,\tfrac1b,\tfrac1c,\tfrac1d)$,
which is at most $\tfrac{101}{105}$, $\tfrac{131}{140}$, $\tfrac{179}{210}$, $\tfrac6b\le1$ according
as $b=3$, $b=4$, $b=5$, $b\ge6$.
\end{proof}

\begin{theorem}[five primes, one multilevel]\label{thm:fiveprimes}
Let $S=T+f$ where $T(n)=\sum_{i=1}^4w_i[q_i\mid n]$ with $w_i\in[0,1]$ and $q_1,\dots,q_4$ powers of
four distinct primes, and $f(n)=\sum_j\alpha_j[p^j\mid n]$ with $\alpha_j\ge0$, $\sum_j\alpha_j\le1$,
for a fifth prime $p$. Then $(H_2)$ holds for $S$, every $m$ and every window. In particular it holds
for five arbitrary primes carrying one arbitrary prime power each.
\end{theorem}
\begin{proof}
For $a\ge0$ and $f\in[0,1]$ the layer identity
$(T+f-a)^+-(T-a)^+=\int_0^1[f>t]\,[T>a-t]\,dt$
holds (both sides equal $(\min(f,1)-(a-T)^+)^+$; the hypothesis $f\le1$ is the per-prime capacity and
is the only one needed). Because $f(n)=A_{v_p(n)}$ with $A_j=\sum_{h\le j}\alpha_h$ nondecreasing and
$A_0=0$, for each $t>0$ the indicator $[f(n)>t]$ is either identically zero or equals $[p^{j(t)}\mid n]$
with $j(t)=\min\{j:A_j>t\}\ge1$; write $d=p^{j(t)}\ge2$. Put $c=1-t$ and let
$Q(t)=\min\{\prod_{i\in B}q_i:\sum_{i\in B}w_i>c\}$ when such a $B$ exists. If $T(n)>c+1$ then
deleting any single $w_i$ leaves a value $>c$, so the set of indices dividing $n$ is feasible for
Lemma~\ref{lem:antichain} and contains a minimal member; hence
$[T>c+1]\le\sum_{A\in\mathcal M_{Q(t)}}[q_A\mid n]$ and, by Lemma~\ref{lem:antichain},
$\#\{k\le M:T(k)>c+1\}\le M\sum_A q_A^{-1}\le M/Q(t)$, an integer bound, so $\le\lfloor M/Q(t)\rfloor$.
Since $d$ is coprime to every $q_i$ we have $T(dk)=T(k)$, so taking $M=\lfloor m/d\rfloor$,
\[
\sum_{k\le m}[d\mid k]\,[T(k)>2-t]\ \le\ \Bigl\lfloor\tfrac m{dQ(t)}\Bigr\rfloor ,
\]
while $[dQ(t)\mid n]\le[d\mid n][T(n)>1-t]$ shows the window contains at least that many points
contributing to the right side. Integrating over $t$ (the integrand is a step function with finitely
many cells) gives
$\sum_{k\le m}\Delta_2(T,f)(k)\le\sum_{b\in I}\Delta_1(T,f)(b)$, and adding
Proposition~\ref{prop:fourprimes} for $T$ gives $(H_2)$ for $S$.
\end{proof}

\begin{remark}[domination after summation, not pointwise]\label{rem:aftersum}
The certificate assembled in the proof --- a finite nonnegative combination of divisor indicators,
namely $\tfrac13\sum_{i<j}(w_i+w_j-1)^+[q_iq_j\mid n]$ plus one term $(v_\ell-u_\ell)[d_\ell Q_\ell\mid n]$
per cell --- dominates $(S-1)^+$ \emph{pointwise} but dominates $(S-2)^+$ only \emph{after summation
over $[1,m]$}. The pointwise statement is false for it: with atoms
$\tfrac{257}{720},\tfrac{421}{720},\tfrac{563}{720},\tfrac{611}{720}$ at $2,3,5,7$ and
$\tfrac{293}{720},\tfrac{427}{720}$ at $11,11^2$, the certificate takes the value $\tfrac{1948}{2160}$
at $n=12705$ while $(S(n)-2)^+=\tfrac{2625}{2160}$. This is not a defect but the mechanism, and it is
what reconciles Theorem~\ref{thm:fiveprimes} with Proposition~\ref{prop:sixobstruction} below.
\end{remark}

\begin{theorem}[any number of primes under a weight cap]\label{thm:cap}
Let $N\ge2$ and $\lambda=\tfrac2{N-1}$. If at most $N$ primes are active and each has total weight at
most $\lambda$, then $(H_2)$ holds.
\end{theorem}
\begin{proof}
Pad to $f_0,\dots,f_{N-1}\in[0,\lambda]$ and set
$G(f)=\sum_{r=0}^{N-1}\max\bigl(0,\min\bigl(\tfrac\lambda N,\min_i(f_i-\lambda a_{ri})\bigr)\bigr)$ with
$a_{ri}=((r+i)\bmod N)/N$. Then $G(f)=\lambda\,\mathbb P(Z_i\le f_i\ \forall i)$ where
$Z_i=\lambda\{U+i/N\}$ and $U$ is uniform on $[0,1)$; the union bound gives
$G\ge\lambda(1-\sum_i(1-f_i/\lambda))=\sum_if_i-2$, and $G\ge0$, so $G\ge(\sum_if_i-2)^+$. On branch
$r$, writing $U=r/N+t$, all the threshold inequalities force $\sum_if_i\ge1+N\lambda t$ because
$\sum_ia_{ri}=\tfrac{N-1}2$, so the admissible $t$-length is at most $(\sum_if_i-1)^+/(N\lambda)$;
multiplying by $\lambda$ and summing the $N$ branches gives $G\le(\sum_if_i-1)^+$. Partitioning each
branch at the finitely many cumulative levels turns $G$ into a nonnegative combination of divisor
indicators, and the interval counts finish as before.
\end{proof}

At five equal half-weights this certificate has value $\tfrac12$, exactly the lower hinge, where the
pair certificate of Proposition~\ref{prop:fourprimes} collapses to $0$.

\begin{proposition}[no pointwise nonnegative certificate at six primes]\label{prop:sixobstruction}
There is no $G(f)=\sum_tc_t\prod_{i:t_i>0}[f_i\ge t_i]$ with all $c_t\ge0$ on $\{0,\tfrac12,1\}^6$
satisfying $(\sum_if_i-2)^+\le G\le(\sum_if_i-1)^+$ everywhere. Equivalently, for the six-prime
system with atoms $\tfrac12$ at $p_i$ and $\tfrac12$ at $p_i^2$, no nonnegative divisor certificate
$C(n)=\sum_d\gamma_d[d\mid n]$ satisfies $(S-2)^+\le C\le(S-1)^+$ pointwise.
\end{proposition}
\begin{proof}
The upper bound at the states of total weight at most $1$ forces the constant, singleton and
$(\tfrac12,\tfrac12)$-pair coefficients to vanish. Let $P,T,H$ be the total coefficient mass of
support $2$, $3$ and $\ge4$. The all-ones state gives $P+T+H\le5$. Each of the twenty states with
exactly three coordinates equal to $1$ has lower hinge $1$, and a threshold vector of support $s$ is
dominated by exactly $\binom{6-s}{3-s}$ of them, so summing gives $4P+T\ge20$. But
$4P+T=4(P+T+H)-3T-4H\le20$, forcing $T=H=0$. Every surviving support-$2$ coefficient has a threshold
equal to $1$ in some coordinate, so $G$ vanishes at the all-half state, where the lower hinge is $1$.
\end{proof}

Proposition~\ref{prop:sixobstruction} is an obstruction to a \emph{method}, not evidence against
$(H_2)$: by Remark~\ref{rem:aftersum} a certificate need only dominate the left side after
summation, and the exhaustive searches reported below found no violation of $(H_2)$ itself. What it
does say is that past five primes any certificate must be signed, or must give up pointwise left
domination, as Theorem~\ref{thm:fiveprimes} does.

\begin{corollary}[the surviving five-prime case]\label{cor:fivesurviving}
If $(H_2)$ fails for a system with at most five active primes, then at least two of those primes carry
more than one nonzero atom, and at least one carries total weight greater than $\tfrac12$.
\end{corollary}
\begin{proof}
Theorem~\ref{thm:fiveprimes} excludes systems with at most one multilevel prime, and
Theorem~\ref{thm:cap} with $N=5$, where $\lambda=\tfrac12$, excludes systems all of whose primes have
total weight at most $\tfrac12$.
\end{proof}

Two remarks delimit these results. The weight cap of Theorem~\ref{thm:cap} is load-bearing and not an
artefact: at $N=6$ with six weights equal to $\tfrac12>\tfrac25=\lambda$ the cyclic certificate falls
below the lower hinge. And the obstruction of Proposition~\ref{prop:sixobstruction} genuinely begins
at six primes: a pointwise nonnegative certificate does exist on $\{0,\tfrac12,1\}^5$. Pairwise
coprimality in Lemma~\ref{lem:antichain} is likewise essential, since $(a,b,c,d)=(2,3,4,5)$ with
$Q=2$ gives $\tfrac{71}{60}>1$.

\begin{remark}[the bridge to unrestricted five primes]\label{rem:K4}
The elimination proof of Theorem~\ref{thm:fiveprimes} needs the antichain bound only for the four
moduli actually present. For an arbitrary multilevel system on four primes $P$ the corresponding
statement is: for every integer $Q\ge2$, the coordinatewise-minimal exponent vectors $e$ with
$D(e)/\max_ip_i^{e_i}\ge Q$, where $D(e)=\prod p_i^{e_i}$, satisfy $Q\sum_e D(e)^{-1}\le1$. This
would give $(H_2)$ for five unrestricted primes. It is open. It holds in every case computed --- all
four-element subsets of the primes below $60$ with $2\le Q\le400$, and $P=\{2,3,5,7\}$ with
$2\le Q\le4000$ --- with maximum $\tfrac{101}{105}$, a margin of under $4\%$. Two structural remarks
make it delicate: the quantity is not monotone in the primes
($K_{\{2,5,7,11\}}(3)=\tfrac{111}{220}<\tfrac{236}{385}=K_{\{3,5,7,11\}}(3)$), and the maximising $Q$
is not always $2$ (for $\{2,5,7,11\}$ it is $Q=4$, with value $\tfrac{37}{55}>\tfrac{213}{385}$), so
the statement does not collapse to Lemma~\ref{lem:antichain}. Even granted, it would not iterate: the
five-prime analogue is false, since
$2\,e_2(\tfrac12,\tfrac13,\tfrac15,\tfrac17,\tfrac1{11})=\tfrac{194}{165}>1$.
\end{remark}

\subsection*{The incidence graph}

The reduction of Sections~\ref{sec:linear}--\ref{sec:h883} passes through the linear programme, and
Proposition~\ref{prop:barrier} caps everything it can give at $4n$. The following bound uses no
relaxation at all, only the count $\#\{b\in I:d\mid b\}\ge\lfloor a_n/d\rfloor$, and is therefore not
subject to that cap. Let $\mathcal P(A)$ be the set of primes dividing $\prod A$, let $H$ be the
bipartite graph on the $n$ inputs and $\mathcal P(A)$ with an edge $(a,p)$ when $p\mid a$, let $c$ be
its number of components and
\[
\beta=\sum_{a\in A}\omega(a)-\bigl|\mathcal P(A)\bigr|+c-n
\]
its cycle rank. (Primes dividing a single input are pendant vertices, so $\beta$ and $c$ are unchanged
if they are deleted; $\beta$ is the cycle rank of the shared-prime subgraph.)

\begin{theorem}[collision--cycle bound]\label{thm:incidence}
$g(A,x)\le n+\beta$ for every $A$ and every $x\ge0$, with no hypothesis on $a_n$.
\end{theorem}
\begin{proof}
Every $a\in A$ satisfies $a\le a_n$, so $\#\{b\in I:a\mid b\}\ge\lfloor a_n/a\rfloor\ge1$; choose any
$b(a)\in I$ divisible by $a$ and let $S=\{b(a):a\in A\}$, a set. The fibres $C_b=\{a:b(a)=b\}$
partition $A$, and distinct fibres have distinct anchors. Write $r_{C,p}(j)=\#\{a\in C:p^j\mid a\}$.
If $r_{C,p}(j)>0$ then $p^j\mid b_C$, so at least $\#\{C:r_{C,p}(j)>0\}$ elements of $S$ are divisible
by $p^j$, whence $c_p(j)-s_p(j)\le\sum_C(r_{C,p}(j)-1)^+\le\sum_C(r_{C,p}(1)-1)^+$ by monotonicity in
$j$; taking the maximum over $j$ gives $\delta_p(S)\le\sum_C(r_{C,p}(1)-1)^+$. Summing over
$p\in\mathcal P(A)$,
\[
\sum_p\delta_p(S)\ \le\ \sum_C\Bigl(\sum_{a\in C}\omega(a)-|\mathcal P(C)|\Bigr)
=\sum_C\bigl(|C|-c_C+\beta_C\bigr),
\]
by Euler's formula in each fibre's own incidence graph $H_C$. Identifying the prime vertices of the
$H_C$ to rebuild $H$ removes one vertex at a time and either merges two components or creates a
cycle, so $\sum_C\beta_C\le\beta$. Hence
$|S|+\sum_p\delta_p(S)\le\sum_C(1+|C|-c_C+\beta_C)\le n+\sum_C\beta_C\le n+\beta$, and
Lemma~\ref{lem:h97-1}'s companion anchor bound $g\le|S|+\sum_p\delta_p(S)$ finishes the proof.
\end{proof}

\begin{corollary}\label{cor:betan}
If $\beta\le n$, equivalently if $\sum_{a\in A}(\omega(a)-2)\le|\mathcal P(A)|-c$, then
$g(A,x)\le2n$.
\end{corollary}

Corollary~\ref{cor:betan} contains the earlier criterion that $g\le2n$ whenever every component of the
incidence graph carries at most one cycle: a pseudoforest has $\beta\le c\le n$. The coefficient of $\beta$ in
Theorem~\ref{thm:incidence} cannot be lowered on the instances below.

\begin{proposition}[sharpness]\label{prop:incsharp}
Let $t\ge2$ and let $w_1<\dots<w_t<q_u<q_v$ be primes all within a factor $\sqrt2$ with
$q_u^2>w_tq_v=:m$. Put $A=\{w_iq_u,w_iq_v:i\le t\}$, the edge set of $K_{2,t}$, so that $n=2t$,
$c=1$ and $\beta=t-1$. If the window of length $m$ can be placed so that $M=w_1\cdots w_tq_uq_v$ is
its only element divisible by two of the demanded primes and neither $q_u^2$ nor $q_v^2$ has a
multiple in it, then $g(A,x)=n+\beta=3t-1$.
\end{proposition}
\begin{proof}
A cover containing $M$ needs $t-1$ further multiples of $q_u$, $t-1$ of $q_v$ and one of each $w_i$,
giving $3t-1$; a cover avoiding $M$ needs at least $t$ multiples of each of $q_u,q_v$ and $t$ more,
giving $3t$.
\end{proof}

Such a placement is not automatic and has to be exhibited. It exists for $t=2,\dots,6$, giving the
exact values $g=5,8,11,14,17$ at $n=4,6,8,10,12$; for instance $t=2$ with
$(w_1,w_2,q_u,q_v)=(59,61,71,73)$ and $x=18\,652\,287$ gives $g=5$, and $t=3$ with
$(67,73,79;97,101)$ and $x=3\,785\,449\,127$ gives $g=8$. In particular $g(10)\ge14$ and
$g(12)\ge17$; the ratio $g/n=\tfrac32-\tfrac1n$ along the family stays below the $(2-o(1))n$ of the
Erd\H{o}s--Sur\'anyi lower bound, as it must, but these are the largest exact values we have
certified at these sizes. Since $\beta=t-1$ grows with $t$, the coefficient of $\beta$ in
Theorem~\ref{thm:incidence} cannot be replaced by any $\theta<1$ at these sizes, and cannot be so
replaced at all if the placement exists for every $t$ --- which we have not proved.

Theorem~\ref{thm:incidence} does not settle the conjecture, and the reason is quantitative rather than
structural: on the family of all $\binom\ell3$ triple products of $\ell$ primes one has
$|\mathcal P(A)|=\ell\asymp n^{1/3}$ and $\beta/n\to2$, so the bound gives only $3n$ there. What is
missing is an arithmetic saving on primitive connected instances whose cycle rank exceeds $n$.

\subsection*{Provenance}
Lemma~\ref{lem:primorial}, Propositions~\ref{prop:fourprimes}--\ref{prop:remains} and Remark~\ref{rem:notight} were
found by GPT-6 (web interface, ``Pro'') in a $41$-minute run on 2026-09-08, whose declared targets --- a proof of the
hinge inequality with threshold $2$, or a refutation at some $c\ge2$ --- were \emph{not} reached.
Lemma~\ref{lem:antichain}, Theorems~\ref{thm:fiveprimes} and~\ref{thm:cap}, Remarks~\ref{rem:aftersum} and~\ref{rem:K4}
and Proposition~\ref{prop:sixobstruction} come from a second $50$-minute run of the same system on the same day, which
likewise did not reach its target: the largest number of primes for which $(H_2)$ is proved without further hypotheses
remains four, and no counterexample was found. That run was refereed by a Claude Opus 5 agent which re-derived the layer
identity, swept the antichain bound over $6.1$ million (quadruple, $Q$) instances, re-verified the cyclic certificate on
$56\,877$ grid points, rebuilt the explicit five-prime certificate over its full period $25\,410$ (reproducing its
pointwise deficit at $n=12705$), and gave an independent exact rational Farkas certificate for
Proposition~\ref{prop:sixobstruction}; it also supplied the two structural obstructions in Remark~\ref{rem:K4}. The
system's own exhaustive audits cover $190\,246\,980$ window instances on five to eight primes without a violation. They were refereed
blind by a Claude Opus 5 agent, which recomputed every constant in exact arithmetic, verified the four-prime
certificate exhaustively over all weight vectors in $\{0,\tfrac16,\dots,1\}^4$ on the primes $2,3,5,7$ for all
$m\le210$ and all windows ($105\,884\,100$ instances), established the optimality of the primorial cutoff and the
sharper bound $m\ge66$, and confirmed that the region of Proposition~\ref{prop:remains} is non-empty. None of the
results of this section is formalised.
The bound of Theorem~\ref{thm:incidence} was found by GPT-6 (web interface, ``Pro'') in a $46$-minute
run on 2026-09-08 aimed at the conjecture itself, which it did not reach; the run also restated three
results of the introduction and of Section~\ref{sec:ksplit} as if new, and those are not repeated here.
Proposition~\ref{prop:incsharp}, the reduction of the run's core-graph corollary to
Corollary~\ref{cor:betan}, and the exact values at $n\le12$ are due to the refereeing Claude Opus 5
agent, which re-derived the proof step by step and recomputed every optimum by an independent dynamic
programme. The orchestrating agent then re-verified the sharpness family at $t=2,3$ and found that the
primes first proposed for it admit no admissible window at $t=2$, which is why the placement is stated
as a hypothesis above and exhibited case by case.
